\begin{figure}[ht]
\centering
\begin{tikzpicture}[x=1cm, y=1cm, font=\small]

% Family-wise error rate (FWER) as a function of the number of
% independent tests m, at nominal per-test alpha = 0.05. Three
% curves: uncorrected (1 - 0.95^m), Bonferroni at alpha/m
% (FWER held at <= 0.05 by construction), and Benjamini-Hochberg
% bound at q = 0.05 (FWER varies with the number of true nulls).

% axes
\draw[->, thick] (0, 0) -- (9.5, 0)
  node[below right, font=\scriptsize] {number of independent tests $m$};
\draw[->, thick] (0, 0) -- (0, 5.5)
  node[left, font=\scriptsize, rotate=90, yshift=0.6cm, anchor=south]
  {Pr(at least one false positive)};

% x ticks: log10 m = 0..3, mapped 0..9 cm
\foreach \mval/\xc in {1/0, 2/0.903, 5/2.097, 10/3, 20/3.903, 50/5.097, 100/6, 500/8.097} {
  \draw (\xc, 0) -- (\xc, -0.1)
    node[below, font=\scriptsize] {\mval};
}
% y ticks
\foreach \pval/\yc in {0/0, 0.2/1, 0.4/2, 0.5/2.5, 0.6/3, 0.8/4, 1.0/5} {
  \draw (0, \yc) -- (-0.1, \yc)
    node[left, font=\scriptsize] {\pval};
}

% Reference at FWER = 0.05
\draw[dashed, gray] (0, 0.25) -- (9.3, 0.25);
\node[font=\scriptsize, gray, anchor=west] at (5.6, 0.55) {nominal $\alpha = 0.05$};

% Uncorrected curve: FWER = 1 - 0.95^m
% In plot coords, x = 3 * log10(m), so m = 10^(x/3)
\draw[thick, engred, smooth]
  plot[domain=0:8.5, samples=80]
  (\x, {5*(1 - pow(0.95, pow(10, \x/3)))});
\node[engred, font=\scriptsize, anchor=west] at (5.5, 4.7)
  {uncorrected};

% Bonferroni: FWER held flat at 0.05 (approximately exact for
% small per-test alpha)
\draw[thick, engblue] (0, 0.25) -- (9.3, 0.25);
\node[engblue, font=\scriptsize, anchor=west] at (3.2, 0.10) {Bonferroni};

% Benjamini-Hochberg FDR control at q=0.05. Under the global null,
% BH controls FWER at q exactly. Under partial null, FWER drifts up
% with the fraction of false hypotheses. We sketch the curve under
% the global null (flat at 0.05) and a partial-null overlay at
% q*(1 + log m / 2) trending up slowly.
\draw[thick, engamber, dashed]
  plot[domain=0:8.5, samples=40]
  (\x, {5*(0.05*(1 + 0.4*\x/3))});
\node[engamber, font=\scriptsize, anchor=south] at (4.5, 0.95)
  {BH (partial null)};

\end{tikzpicture}
\caption[FWER vs number of tests under three corrections]{Family-
wise error rate (probability of at least one false positive) as a
function of the number of independent tests $m$, each at nominal
$\alpha = 0.05$. With \emph{no} correction (engineering red), the
FWER rises rapidly: by $m = 14$ it is above $0.5$, and by $m = 50$
it is above $0.92$. \emph{Bonferroni} correction (engineering blue)
holds FWER at $0.05$ exactly by lowering each per-test threshold to
$\alpha/m$. \emph{Benjamini-Hochberg} false-discovery-rate control
at $q = 0.05$ (engineering amber, dashed) keeps the expected
proportion of false discoveries among rejections at $q$; under a
partial null (most hypotheses true, some false) the resulting FWER
drifts upward slowly with $m$, the cost of trading a stricter FWER
guarantee for power on the real effects. The $x$-axis is
logarithmic.}
\label{fig:vol01:ch08:bonferroni-vs-fdr}
\end{figure}
